\subsection{Definition of prime ideals}

DEFINITION 1. An ideal p of a ring A is called prime if the ring A/p is an integral domain.

By this definition, an ideal p of a ring A is prime if the following two conditions hold:

\begin{enumerate}
\def\labelenumi{(\arabic{enumi})}
\item
  \(\mathfrak{p} \neq \mathbf{A}\) ;
\item
  if \(x, y\) are two elements of \(A\) such that \(x \notin p\) and \(y \notin p\) , then \(xy \notin p\) .
\end{enumerate}

These conditions can also be expressed by saying that the product of any finite family of elements of \(\mathbb{C}\mathfrak{p}\) belongs to \(\mathbb{C}\mathfrak{p}\) , as applying this condition to the empty set yields \(1\notin\mathfrak{p}\) .

A maximal ideal m of A is prime since A/m is a field; then it follows from Krull's theorem (Algebra, Chapter I, § 8, no. 7, Theorem 2) that every ideal of

A other than A is contained in at least one prime ideal. In particular, for prime ideals to exist in a ring A, it is necessary and sufficient that A be not reduced to 0.

Let \(f: \mathbf{A} \to \mathbf{B}\) be a ring homomorphism and \(q\) an ideal of \(\mathbf{B}\) . Set \(p = \bar{f}^{-1}(q)\) ; the homomorphism \(\overline{f}: \mathbf{A}/p \to \mathbf{B}/q\) derived from \(f\) by taking quotients is injective. Suppose that \(q\) is prime; as the ring \(B/q\) is an integral domain, so is \(A/p\) , being isomorphic to a subring of \(B/q\) ; consequently the ideal \(p = \bar{f}^{-1}(q)\) is prime. In particular, let \(A\) be a subring of \(B\) ; for every ideal \(q\) of \(B, q \cap A\) is a prime ideal \(A\) . If \(f\) is surjective, \(\overline{f}\) is an isomorphism; the conditions `` \(p\) is prime'' and `` \(q\) is prime'' are then equivalent. Hence, if \(p\) and \(a\) are ideals of \(A\) such that \(a \subset p\) , a necessary and sufficient condition for \(p\) to be prime is that \(p/a\) be prime in \(A/a\) .

PROPOSITION 1. Let A be a ring, \(\mathfrak{a}_1, \mathfrak{a}_2, \ldots, \mathfrak{a}_n\) ideals of A and p a prime ideal of A. If p contains the product \(\mathfrak{a}_1\mathfrak{a}_2 \ldots \mathfrak{a}_n\) , it contains at least one of the \(\mathfrak{a}_t\) .

Suppose in fact that \(\mathfrak{p}\) contains none of the \(a_{i}\) . For \(1 \leqslant i \leqslant n\) there exists then an element \(s_{i} \in a_{i} \cap \mathbb{C}p\) ; then \(s = s_{1}s_{2} \ldots s_{n}\) is contained in \(a_{1}a_{2} \ldots a_{n}\) and is not contained in \(\mathfrak{p}\) , which is absurd.

COROLLARY. Let \(\mathfrak{m}\) be a maximal ideal of \(\mathbf{A}\) ; for every integer \(n > 0\) , the only prime ideal containing \(\mathfrak{m}^n\) is \(\mathfrak{m}\) .

Such an ideal \(\mathfrak{p}\) must contain \(\mathfrak{m}\) by Proposition 1 applied to \(a_{i} = \mathfrak{m}\) for \(1 \leqslant i \leqslant n\) ; as \(\mathfrak{m}\) is maximal, \(\mathfrak{p} = \mathfrak{m}\) .

PROPOSITION 2. Let A be a ring, a a non-empty set of A which is closed under addition and multiplication and \((\mathfrak{p}_{i})_{i\in I}\) a non-empty finite family of ideals of A. Suppose that a is contained in the union of the \(p_{i}\) and that at most two of the \(p_{i}\) are not prime. Then a is contained in one of the \(p_{i}\) .

We argue by induction on \(n = \operatorname{Card}(\mathbf{I})\) ; the proposition is trivial if \(n = 1\) . Suppose that \(n \geqslant 2\) ; if there exists an index \(j\) such that \(\mathfrak{a} \cap \mathfrak{p}_j \subset \bigcup_{i \neq j} \mathfrak{p}_i\) , the set \(\mathfrak{a}\) , which is the union of the \(\mathfrak{a} \cap \mathfrak{p}_i\) where \(i \in \mathbf{I}\) , is contained in \(\bigcup_{i \neq j} \mathfrak{p}_i\) and hence in one of the \(\mathfrak{p}_i\) by the induction hypothesis. Suppose then that such an index does not exist; for every \(j \in \mathbf{I}\) let \(y_j\) be an element of \(\mathfrak{a} \cap \mathfrak{p}_j\) not belonging to any \(\mathfrak{p}_i\) for \(i \neq j\) . Let \(k\) be an element of \(\mathbf{I}\) chosen in such a way that \(\mathfrak{p}_k\) is prime if \(n > 2\) and chosen arbitrarily if \(n = 2\) ; let \(z = y_k + \prod_{i \neq k} y_i\) . Then \(z \in \mathfrak{a}\) , since \(\mathfrak{a}\) is closed under addition and multiplication; if \(j \neq k, \prod_{i \neq k} y_i\) belongs to \(\mathfrak{p}_j\) , but \(y_k \notin \mathfrak{p}_j\) , whence \(z \notin \mathfrak{p}_j\) . On the other hand, \(\prod_{i \neq k} y_i\) does not belong to \(\mathfrak{p}_k\) , as none of the factors \(y_i (i \neq k)\) belongs to it and \(\mathfrak{p}_k\) is prime if \(n - 1 > 1\) ; as \(y_k \in \mathfrak{p}_k\) , \(z\) does not belong to \(\mathfrak{p}_k\) and the proposition is established.

